%% =====================================================================
%%  structures/legislativo/structure.tex  --  Forma parlamentaria de los
%%  productos legislativos (proyecto de ley, dictamen, mocion).
%% ---------------------------------------------------------------------
%%  Método   : la ESTRUCTURA (que componentes, en que orden) la fija el
%%             esquema del AWF (esquemas/proyecto-de-ley.tex, dictamen.tex,
%%             mocion.tex; Reglamento arts. 68, 70, 75; Manual TL 2021).
%%             Aqui solo va la FORMA: como se ve cada pieza, segun el
%%             Manual de Tecnica Legislativa (calibre-id 10098): el titulo
%%             de la ley centrado; "EXPOSICION DE MOTIVOS" centrada, en
%%             mayuscula y negrita (sec. X.c.4); los articulos "Articulo N.
%%             Epigrafe"; las disposiciones complementarias por bloque.
%%  Uso      : \componente{clave} abre cada pieza; la cabecera la resuelve
%%             \esq@cabecera segun el nivel, salvo las piezas MUDAS
%%             (presentacion, titulo, firmas, sumilla, encabezado, estilo),
%%             que no llevan rotulo y se escriben con las macros de abajo.
%% =====================================================================
\ifcsname despachoEstructuraLegislativo\endcsname
  \expandafter\endinput
\fi
\expandafter\def\csname despachoEstructuraLegislativo\endcsname{cargada}
\RequirePackage{needspace}

%% ── Rotulos ─────────────────────────────────────────────────────────
%% Nivel 1 (formula legal, exposicion de motivos, analisis...): centrado,
%% mayuscula, negrita. Nivel 2 (fundamentos, antecedentes, costo-beneficio,
%% analisis tecnico...): numerado en romano, como las secciones del informe.
%% Espacios compactos (2026-10-01, revisión del proyecto real): el entorno
%% `center` sumaba su topsep a cada \vspace y dejaba la fórmula legal entre
%% dos franjas de blanco; ahora el centrado es un párrafo y el aire, medido.
\newcounter{seccionlegislativa}
\newcommand{\rotulolegislativo}[1]{%
  \par\needspace{5\baselineskip}%
  \vspace{0.5\baselineskip}%
  {\centering\fuentedisplaysemi\color{colortitulos}\MakeUppercase{#1}\par}%
  \vspace{0.15\baselineskip}\setcounter{seccionlegislativa}{0}}
%% \seccionlegislativa[II]{Efecto de la vigencia...}: el numero romano lo da
%% el TEXTO cuando lo trae (el generador lo toma de sus apartados: «III.1» ->
%% III); sin argumento opcional, el contador. El titulo va en mayusculas y
%% minusculas: la jerarquia la dan la fuente, el color y el numero, no las
%% mayusculas (2026-10-01). Si el titulo ocupa dos lineas, la segunda se
%% alinea con el texto de la primera, no con el numero.
\newcommand{\seccionlegislativa}[2][]{%
  \par\needspace{5\baselineskip}%
  \stepcounter{seccionlegislativa}%
  \vspace{0.6\baselineskip}%
  \sbox0{\fuentedisplaysemi\ifstrempty{#1}{\Roman{seccionlegislativa}}{#1}.\ }%
  \noindent\hangindent=\wd0\hangafter=1
  {\fuentedisplaysemi\color{colortitulos}\usebox0 #2}\par
  \nopagebreak\vspace{0.1\baselineskip}}

%% Las piezas MUDAS: el esquema las exige y las registra, pero no llevan
%% rotulo; su forma la ponen las macros de mas abajo.
%% (Se guarda y restaura el catcode de `@`: esta estructura la carga la CLASE,
%% donde `@` es letra, y un \makeatother suelto rompe lo que la clase haga despues.)
\edef\legislativoCatcode{\catcode`\noexpand\@=\the\catcode`\@\relax}
\makeatletter
\newcommand{\legislativo@muda}[1]{\expandafter\def\csname legmuda@#1\endcsname{}}
\legislativo@muda{presentacion}\legislativo@muda{titulo}\legislativo@muda{firmas}
\legislativo@muda{sumilla}\legislativo@muda{encabezado}\legislativo@muda{estilo}
\legislativo@muda{firma}\legislativo@muda{destinatario}
%% La formula legal y sus partes tampoco llevan rotulo del esquema: una ley
%% no dice "Parte sustantiva" (Manual sec. VI); las disposiciones
%% complementarias abren su bloque con \bloquedisposiciones, y el titulo
%% preliminar o el "TEXTO SUSTITUTORIO" del dictamen, con \rotulolegislativo.
\legislativo@muda{formula}\legislativo@muda{preliminar}\legislativo@muda{sustantiva}
\legislativo@muda{dcf}\legislativo@muda{dct}\legislativo@muda{dcm}\legislativo@muda{dcd}
\legislativo@muda{anexo-formula}
%% La EXPOSICION DE MOTIVOS empieza siempre en hoja nueva (2026-10-01).
%% \cabeceraexplicita{...}: el generador decide como se rotula la PROXIMA
%% pieza (el numero real de la seccion, o un apartado «I.17.» cuando la
%% pieza vive dentro de otra seccion); sin ella, el rotulo del esquema.
\newcommand{\cabeceraexplicita}[1]{\gdef\legcabeceraexplicita{#1}}
\renewcommand{\esq@cabecera}[2]{% {titulo}{clave}
  % La clave llega como macro sin expandir: se expande antes de comparar (con
  % \ifstrequal a secas la comparacion fallaba y no habia salto de pagina).
  \edef\legclaveactual{#2}\ifdefstring{\legclaveactual}{exposicion}{\clearpage}{}%
  \ifdefined\legcabeceraexplicita
    \legcabeceraexplicita\global\let\legcabeceraexplicita\undefined
  \else
    \ifcsname legmuda@#2\endcsname\else
      \ifdefstring{\esq@nivelactual}{2}{\seccionlegislativa{#1}}{\rotulolegislativo{#1}}%
    \fi
  \fi}
%% \proyectoenavance: un proyecto EN ELABORACION (primer avance sin exposicion
%% de motivos, 2026-10-01) se compone aunque falten componentes obligatorios
%% del esquema: la verificacion avisa en el log en vez de detener la
%% compilacion. Lo pide el generador del SGDP solo cuando faltan partes, y la
%% interfaz dice cuales; la validacion completa sigue siendo la de siempre.
\newcommand{\proyectoenavance}{%
  \renewcommand{\esq@verifica}{\typeout{[ESQUEMA] \esq@nombre: proyecto en avance; se compone sin las partes obligatorias que faltan.}}}
\legislativoCatcode

%% ── La formula legal ─────────────────────────────────────────────────
%% \titulodelaley{Ley que ...}   (Manual sec. VI.c.3: el titulo de la ley)
%% Justificado con la ULTIMA linea centrada (como el Word de la Camara,
%% 2026-10-01: «no es centrado, es justificado»): el truco de los tres
%% pegamentos reparte el blanco a los lados solo en la ultima linea.
\newcommand{\titulodelaley}[1]{%
  \par\vspace{0.1\baselineskip}%
  {\leftskip=0pt plus 1fil\rightskip=0pt plus -1fil\parfillskip=0pt plus 2fil
   \fuentedisplaysemi\MakeUppercase{#1}\par}%
  \vspace{0.45\baselineskip}}
%% \articulo{1}{Objeto de la ley}  -> "Articulo 1. Objeto de la ley" y el texto
%% El separador entre numero y epigrafe es configurable (2026-10-01): el Manual
%% (2021) dice «Articulo 1.», pero la Camara de Diputados usa «Articulo 1.-»
%% (\setseparadorarticulo{.-\ }). Lo mismo para las disposiciones
%% («PRIMERA. -»: \setseparadordisposicion{. -\ }).
\newcommand{\despachoSepArticulo}{.\ }
\newcommand{\setseparadorarticulo}[1]{\renewcommand{\despachoSepArticulo}{#1}}
\newcommand{\despachoSepDisposicion}{.\ }
\newcommand{\setseparadordisposicion}[1]{\renewcommand{\despachoSepDisposicion}{#1}}
\newcommand{\articulo}[2]{%
  \par\needspace{3\baselineskip}%
  \vspace{0.45\baselineskip}%
  \noindent{\seriesemi Art\'iculo #1\despachoSepArticulo#2}\par\nopagebreak\noindent\ignorespaces}
%% \divisionlegislativa{T\'ITULO I}{Disposiciones generales}: las divisiones de
%% la formula legal (titulo, capitulo, seccion; Manual de Tecnica Legislativa
%% sec. VI): centradas, en mayuscula y negrita, la denominacion debajo
%% (segunda proposicion del despacho, 2026-10-01). Sin denominacion, solo
%% el rotulo.
\newcommand{\divisionlegislativa}[2]{%
  \par\needspace{6\baselineskip}%
  \vspace{0.8\baselineskip}%
  {\centering\fuentedisplaysemi\MakeUppercase{#1}\par
   \ifstrempty{#2}{}{\vspace{0.15\baselineskip}\MakeUppercase{#2}\par}}%
  \nopagebreak\vspace{0.2\baselineskip}}
%% \bloquedisposiciones{FINALES} -> "DISPOSICIONES COMPLEMENTARIAS FINALES"
\newcommand{\bloquedisposiciones}[1]{%
  \par\needspace{5\baselineskip}%
  \vspace{0.8\baselineskip}%
  {\centering\fuentedisplaysemi DISPOSICIONES COMPLEMENTARIAS \MakeUppercase{#1}\par}%
  \vspace{0.1\baselineskip}}
%% \bloquedisposicion{transitoria} -> "DISPOSICION COMPLEMENTARIA TRANSITORIA"
%% (singular, cuando el bloque tiene una sola; 2026-10-01)
\newcommand{\bloquedisposicion}[1]{%
  \par\needspace{5\baselineskip}%
  \vspace{0.8\baselineskip}%
  {\centering\fuentedisplaysemi DISPOSICI\'ON COMPLEMENTARIA \MakeUppercase{#1}\par}%
  \vspace{0.1\baselineskip}}
%% \disposicion{PRIMERA}{Vigencia} -> "PRIMERA. Vigencia" y el texto
\newcommand{\disposicion}[2]{%
  \par\needspace{3\baselineskip}%
  \vspace{0.45\baselineskip}%
  \noindent{\seriesemi\MakeUppercase{#1}\despachoSepDisposicion#2}\par\nopagebreak\noindent\ignorespaces}

%% ── Piezas del proyecto de la Camara de Diputados (2026-10-01) ───────
%% Las trae el primer proyecto real del despacho (piloto de la plataforma).
%% \proposicionlegislativa{Que previene...} -> rotulo en mayuscula y negrita al
%% comienzo de la primera hoja, como en el proyecto presentado: un BLOQUE
%% justificado que ocupa los tres quintos derechos de la caja (a su izquierda
%% queda el lugar del sello de firma digital). Antes era un texto alineado a
%% la derecha con el borde izquierdo irregular (revision 2026-10-01).
\newcommand{\proposicionlegislativa}[1]{%
  \par\noindent\hfill
  \begin{minipage}[t]{0.6\textwidth}%
    \setlength{\parfillskip}{0pt plus 1fil}\setlength{\emergencystretch}{1.5em}%
    {\fuentedisplaysemi\MakeUppercase{Proposici\'on legislativa #1}\par}%
  \end{minipage}\par\vspace{0.8\baselineskip}}
%% \apartadolegislativo[I.10.1.]{Precision del estandar...} -> subtitulo de la
%% exposicion de motivos con su numeracion decimal. El numero va en su caja y
%% el titulo largo se alinea con su primera palabra: numero -> titulo ->
%% contenido, sin heredar la sangria de los parrafos (2026-10-01).
\newcommand{\apartadolegislativo}[2][]{%
  \par\needspace{3\baselineskip}%
  \vspace{0.4\baselineskip}%
  \sbox0{\seriesemi\ifstrempty{#1}{}{#1\ }}%
  \noindent\hangindent=\wd0\hangafter=1
  {\seriesemi\usebox0 #2}\par\nopagebreak}

%% \itemlegislativo{a)}{texto}: un inciso (a), b), 1., i), -) con SANGRIA
%% COMPLETA: la marca al margen y TODAS las lineas del texto alineadas con
%% la primera palabra (pedido 2026-10-01). Es una lista de un solo elemento
%% (no \hangindent): asi se anida dentro del texto citado, cuyas listas usan
%% \parshape y anularian la sangria colgante. El ancho minimo iguala los
%% incisos cortos de una misma serie.
\newlength{\legsangriaincisos}
\setlength{\legsangriaincisos}{1.7em}
\newlength{\legmarcainciso}
\newcommand{\itemlegislativo}[2]{%
  \par\settowidth{\legmarcainciso}{#1\ }%
  \ifdim\legmarcainciso<\legsangriaincisos\setlength{\legmarcainciso}{\legsangriaincisos}\fi
  \begin{list}{}{\setlength{\leftmargin}{\legmarcainciso}\setlength{\labelwidth}{\legmarcainciso}%
    \setlength{\labelsep}{0pt}\setlength{\itemindent}{0pt}\setlength{\listparindent}{0pt}%
    \setlength{\topsep}{0pt}\setlength{\partopsep}{0pt}\setlength{\itemsep}{0pt}%
    \setlength{\parsep}{\parskip}\renewcommand{\makelabel}[1]{##1\hfil}}%
  \item[#1]#2%
  \end{list}}
%% \firmadigital{Nombre}{Cargo}: el LUGAR de la firma digital del autor o de un
%% adherente, en cuadricula de dos columnas (\cuadrofirmas). Solo el espacio:
%% sin nombre, cargo, linea ni rotulo (2026-10-01); el sello de la firma
%% digital trae los suyos. Nombre y cargo siguen en la llamada para saber
%% cuantos espacios hacen falta y para la bitacora del .tex.
\newlength{\alturafirmadigital}
\setlength{\alturafirmadigital}{2.4cm}
\newcounter{firmadigital}
\newcommand{\firmadigital}[2]{%
  \stepcounter{firmadigital}%
  \begin{minipage}[t][\alturafirmadigital]{0.47\textwidth}\mbox{}\end{minipage}%
  \ifodd\value{firmadigital}\hfill\else\par\noindent\fi}
\newenvironment{cuadrofirmas}{%
  \par\needspace{\alturafirmadigital}\vspace{0.8\baselineskip}\setcounter{firmadigital}{0}\noindent}{\par}

%% textocitado: el texto que el proyecto cita de otra norma («en los
%% siguientes terminos: “Articulo 359-D.- …”»), sangrado a la izquierda como
%% en el proyecto presentado (2026-10-01). Dentro, \rubrocitado{} da el
%% epigrafe del articulo citado sin confundirlo con un \articulo propio.
\newenvironment{textocitado}{%
  \begin{list}{}{\setlength{\leftmargin}{1.1cm}\setlength{\rightmargin}{0pt}%
    \setlength{\topsep}{0.35\baselineskip}\setlength{\parsep}{0.35\baselineskip}%
    \setlength{\itemsep}{0pt}\setlength{\listparindent}{0pt}}\item\relax}%
  {\end{list}}
\newcommand{\rubrocitado}[1]{%
  \needspace{3\baselineskip}%
  {\seriesemi #1}\par\nopagebreak}

%% ── Piezas mudas ─────────────────────────────────────────────────────
%% \sumilla{...}  (dictamen, mocion: la sumilla va primero, en cursiva)
\newcommand{\sumilla}[1]{%
  \noindent{\seriesemi Sumilla:}\ \textit{#1}\par\vspace{0.6\baselineskip}}
%% \firmaslegislativas -> el bloque de firma del titular de la iniciativa
\newcommand{\firmaslegislativas}{%
  \vspace{2\baselineskip}%
  \begin{flushright}
    \rule{6.5cm}{0.4pt}\\
    {\seriesemi\meta{firmante}}\\
    \meta{cargofirmante}
  \end{flushright}}
%% \formuladeestilo -> el cierre del dictamen (Manual sec. XI.c.10)
\newcommand{\formuladeestilo}[1][Dese cuenta.\\ Sala de la Comisi\'on.]{%
  \vspace{1\baselineskip}\noindent #1\par
  \noindent\meta{lugar}, \ifmetadef{fecha}{\meta{fecha}}{\today}\par}
